\section{Conclusion}

\subsection{Hypothesis 1}

The obtained results vary completely from the expected values as although velocity is supposed to increase as channel efficiency increases, the results show the very opposite happening at all sites except for at Site 8, where the data corresponds to one another. This goes against Bradshaw's model where velocity increases as we move downstream.



\subsection{Hypothesis 2}

As mentioned in the hypothesis, increase in the river's discharge does correlate to the channel width and channel depth to a certain extent. When comparing the discharge with the channel width, one can notice a clear trend with the river discharge increasing as the width increased and the discharge decreasing as the channel width decreased. The only anomaly is at Site 8, where the channel width increases to 4.24 metres in comparison to Site 7 at 3.87 metres, however, the discharge decreases from 4.31 cubic metres per second. Channel width corresponds to Bradshaw's Model where it increases as we move downstream, however, channel depth doesn't support this. Discharge also supports Bradshaw's Model as it shows an increase moving downstream.




\subsection{Hypothesis 3}

Results obtained for the Bedload size completely support the hypothesis that they decrease in size as we move downstream, except for at Site 6, where it increases to 43.9mm from 38.3mm before dropping to 23.0mm at Site 7. More importantly, these results wholly support the Bradshaw Model.